\begin{textsample}{POS Dim 1 – vem\_ed\_11 – Score 13.00 – vem\_ed\_11/t048.txt}  \label{ex:f1_pos_147}
PCIs com Índia valorizam diversidade cultural Em 2021, a equipe dos PCIs / Cesu \textbf{iniciou} colaborações com o Symbiosis Centre for Management Studies, localizado em Pune, oeste da Índia.
Os principais desafios contornados \textbf{foram} o fuso \textbf{horário}, com 8 horas e meia de \textbf{diferença}, e o tamanho das \textbf{turmas} indianas – existem salas com mais de 100 alunos.
Lark \textbf{foi} a principal ferramenta de comunicação \textbf{utilizada}, para realização dos \textbf{encontros} \textbf{síncronos} e compartilhamento de informações.
Além disso, utilizou-se bastante o WhatsApp para \textbf{interação} entre os participantes.
Dois projetos \textbf{foram} realizados no segundo semestre de 2021: um sobre satisfação no trabalho, nas Fatecs Americana (Ricardo Bertoni Pompeu), Sumaré (Danilo Sergio Sorroce e Rafaeli Cardozo Modolo Begalli) e, na instituição indiana, a professora Nehajoan Panackal.
Nesse projeto participaram cerca de 50 alunos de \textbf{cada} país.
Situações-problema sobre satisfação no trabalho \textbf{eram} apresentadas em vídeo e os grupos propunham soluções. “Fomos surpreendidos pela qualidade do trabalho \textbf{que} encontramos com esse novo parceiro”, comenta Ana Carolina Freschi, da equipe dos PCIs, \textbf{que} acompanhou o projeto. “Tudo correu muito bem, Neha \textbf{é} uma professora muito engajada; os alunos tiveram uma \textbf{interação} muito grande.
Os indianos \textbf{foram} extremamente generosos”, comemora Pompeu.
O outro projeto, sobre \textbf{impacto} da Covid-19 em padrões de investimentos pessoais, \textbf{foi} conduzido pelos professores Ricardo Trovão (Fatec Itaquaquecetuba) e Shreya Virani (Symbiosis).
Os PCIs \textbf{foram} muito \textbf{bem-sucedidos} e bem avaliados, tanto pelos docentes brasileiros \textbf{quanto} pelas professoras indianas.
Portanto, as colaborações continuam em 2022 e se estendem para outras áreas como marketing.
Para além dos \textbf{conteúdos} técnicos, houve o enriquecimento cultural dos estudantes.
Durante a realização do PCI, em 4 de novembro, ocorreu o Diwali – festival hindu \textbf{que} celebra as divindades femininas Kali, Lakshmi e Sarasvati, culminando com a vitória das luzes sobre as trevas.
Os estudantes indianos se preocuparam em preparar uma apresentação sobre o Diwali para os brasileiros, \textbf{contando} sobre comidas, vestimentas, músicas. “Em poucas semanas, brasileiros e indianos conseguiram resolver problemas, interagir, trocar informações. \textbf{Foi} um sucesso”, afirma Freschi.

% matched lemmas: bem-sucedido, cada, contar, conteúdo, diferença, encontro, horário, impacto, iniciar, interação, quanto, que, ser, síncrono, turma, utilizar
\end{textsample}
